\documentclass[11pt]{article}
\usepackage[margin=1in]{geometry}
\usepackage{times}
\usepackage{hyperref}
\hypersetup{colorlinks=true, linkcolor=black, urlcolor=blue, citecolor=black}
\title{A Filed Patent Application Is a Claim, Not a Confirmation: The Pais Patents on Inertial Mass Reduction}
\author{Flyxion \\ \small Independent Researcher}
\date{}
\begin{document}
\maketitle

\begin{abstract}
Salvatore Cezar Pais, while employed as an engineer for the United States Navy, filed and was granted several patents between roughly 2017 and 2021 describing a "high frequency gravitational wave generator," a "craft using an inertial mass reduction device," and a room-temperature superconductor, with the inertial mass reduction patent specifically describing a mechanism for locally reducing an object's effective inertial mass via a vibrating, high-energy-density electromagnetic cavity. This paper argues that the Navy's own internal justification for granting and prioritizing patent examination of these applications, citing competitive pressure from other nations rather than internal experimental confirmation, has been publicly reported and directly contradicts the popular reading of the patents as evidence of a demonstrated working technology, that Navy officials who reviewed the underlying claims, including in sworn testimony, have described the inventions as unproven and not independently verified in any operational or laboratory demonstration, and that the patent grant process, as addressed in the earlier patent-office entry in this series, does not certify that a described invention functions as claimed, a limitation that applies to these patents as directly as to any of the antigravity patents surveyed there.
\end{abstract}

\section{Introduction}

Salvatore Cezar Pais, an engineer at what was then the Naval Air Warfare Center Aircraft Division, filed and was granted United States patents describing several exotic technologies, including a "craft using an inertial mass reduction device" (granted 2018), a "high frequency gravitational wave generator" (granted 2019), and a room-temperature superconductor (granted 2019), each accompanied by patent text invoking a mix of established and speculative physics to describe mechanisms including vacuum energy manipulation, macroscopic quantum coherence effects, and localized manipulation of what the patents describe as an object's effective inertial mass, generating substantial public and media attention, particularly amid concurrent public discussion of unidentified aerial phenomena, as apparent evidence of an advanced, functioning Navy antigravity or exotic propulsion program.

\section{What Reporting on the Patents' Internal Navy History Has Established}

Subsequent investigative reporting, including Freedom of Information Act document releases, revealed that Pais's patent applications initially faced examiner rejection on grounds including a lack of enablement, meaning the application did not adequately describe how to actually build a functioning device, and that Navy officials intervened to advocate for the patents' approval specifically by citing the strategic value of the United States holding the patent rights should the technology prove implementable, including internal correspondence explicitly stating that Navy leadership could not confirm the described inertial mass reduction effect had been experimentally demonstrated, but supported patenting the concept partly to prevent a strategic competitor, reportedly China, from patenting or developing the technology first, a rationale addressing patent strategy and technological competition rather than confirming the underlying physics.

\section{Why This Directly Undercuts the "Working Technology" Reading}

A patent applicant and the reviewing agency explicitly stating, in internal correspondence later made public, that the described effect has not been experimentally confirmed is a direct, first-party contradiction of the popular narrative that these patents represent evidence of a demonstrated, functioning Navy antigravity technology; the patents describe theoretical mechanisms Pais proposed and the Navy chose to secure patent rights over as a strategic hedge, not confirmed experimental results, a distinction the Navy's own internal documentation draws explicitly and that patent grant, as addressed in the general patent-office entry elsewhere in this series, does not itself resolve, since patent examiners assess novelty and, following the enablement objection's eventual resolution in Pais's case, a sufficiently detailed description for a skilled practitioner to attempt construction, not whether the invention actually works as claimed.

\section{The Physics the Patents Themselves Would Need to Overcome}

The inertial mass reduction patent's proposed mechanism, vibrating a cavity at high frequency to create a localized manipulation of the quantum vacuum sufficient to alter an object's inertial mass, faces the same scale-mismatch objection raised against the zero-point energy antigravity hybrid claims addressed elsewhere in this series: the vacuum energy density accessible to any laboratory-scale electromagnetic cavity, however cleverly resonated, remains many orders of magnitude short of what would be required to produce a macroscopically significant inertial mass change, a gap the patent text does not close with a quantitative derivation showing the claimed effect size is achievable at the described cavity dimensions and power levels, leaving the application, like the zero-point energy entry addressed elsewhere in this series, without a demonstrated bridge across the relevant scale gap.

\section{The Entropic Gravity Framing}

The high frequency gravitational wave generator patent specifically proposes generating gravitational radiation at power levels and frequencies vastly beyond anything the general theory of relativity's own well confirmed predictions, independently validated by LIGO's detection of gravitational waves from astrophysical-scale events, would suggest achievable by any laboratory apparatus; on the entropic account of gravitation addressed throughout this series, as on the standard general-relativistic account, generating meaningful gravitational radiation requires accelerating mass-energy configurations at a scale entirely disproportionate to anything a benchtop or vehicle-scale cavity resonator could produce, the same objection raised against the Woodward-Mach effect thruster's more modest and more carefully derived claims elsewhere in this series, applied here to a considerably less quantitatively developed patent description.

\section{Objections}

Supporters note that Pais's patents were reviewed and approved by professional patent examiners at the United States Patent and Trademark Office, and argue this constitutes a meaningful form of technical vetting beyond what less formally examined claims elsewhere in this series have received. As the patent-office entry elsewhere in this series addresses in detail, patent examination assesses novelty and adequate description rather than functional verification, and the specific, publicly documented history of Pais's applications, including the initial enablement rejection and the internal Navy correspondence explicitly declining to confirm experimental verification while advocating for approval on strategic patent-denial grounds, is an unusually well documented instance of exactly this distinction being drawn by the reviewing parties themselves.

\section{Conclusion}

The Navy's own internal correspondence regarding the Pais patents explicitly declines to confirm experimental verification of the claimed effects, framing the patent applications as a strategic hedge against foreign patent claims rather than as documentation of a demonstrated technology, and the underlying physics described, particularly regarding vacuum energy manipulation and gravitational wave generation at laboratory scale, faces the same scale-mismatch objections raised against comparable claims addressed in detail elsewhere in this series, leaving the patents as a notable case study in how patent grant is popularly misread as functional confirmation considerably more clearly than most other entries surveyed here, given how directly the Navy's own documented internal position contradicts that reading.

\end{document}
